\begin{table}[t]
\centering
\small
\caption{Comparison with published approaches for automatically identifying
security weaknesses in LLM tools and MCP servers. Existing approaches either
target predefined vulnerability classes or specific agent--tool attacks.
\project{} instead derives argument-specific security properties from tool
metadata and uses them to dynamically discover semantic implementation bugs.}
\label{tab:motivation}
\resizebox{\textwidth}{!}{
\begin{tabular}{l|c|c|c|c|c|c}
\toprule
\textbf{System} &
\textbf{Tool Impl.} &
\textbf{Dynamic} &
\textbf{Argument Semantics} &
\textbf{Broad Bug Classes} &
\textbf{Semantic Oracle} &
\textbf{Auto. Harness} \\
\midrule

MCP at First Glance~\cite{mcp-first-glance} &
\cmark &
\xmark &
\xmark &
\partialmark &
\xmark &
N/A \\

MCP-SecLint~\cite{mcp-seclint} &
\cmark &
\xmark &
\xmark &
\xmark &
\xmark &
N/A \\

Cisco MCP Scanner~\cite{cisco2025mcpscanner} &
\partialmark &
\partialmark &
\xmark &
\partialmark &
\xmark &
\cmark \\

MCPSecBench~\cite{yang2025mcpsecbench} &
\partialmark &
\cmark &
\xmark &
\partialmark &
\xmark &
\xmark \\

MCP-Scanner~\cite{mcp-scanner} &
\partialmark &
\cmark &
\xmark &
\partialmark &
\xmark &
\cmark \\

% MCPSafetyScanner~\cite{mcp-safety-audit} &
% \partialmark &
% \cmark &
% \partialmark &
% \partialmark &
% \xmark &
% \cmark \\

Chord~\cite{les-dissonances} &
\xmark &
\cmark &
\partialmark &
\xmark &
\partialmark &
\cmark \\

VIPER-MCP~\cite{viper-mcp} &
\cmark &
\cmark &
\partialmark &
\partialmark &
\xmark &
\cmark \\

MCPZoo~\cite{mcpzoo} &
\partialmark &
\partialmark &
\xmark &
\xmark &
\xmark &
\cmark \\

Description-code inconsistency~\cite{li2026don,shi2026description} &
\cmark &
\partialmark &
\partialmark &
\xmark &
\partialmark &
\cmark \\

\midrule
\textbf{\project{}} &
\cmark &
\cmark &
\cmark &
\cmark &
\cmark &
\cmark \\
\bottomrule
\end{tabular}
}
\end{table}